\section{多模态视频理解：声音、视觉和融合}

\subsection{为什么要把音频也算进来}

视频并不只有视觉。很多场景中，音频同样提供重要线索：

\begin{itemize}
    \item 人在说话时，声音能帮助识别是谁在说。
    \item 乐器演奏时，画面和声音可以互相约束。
    \item 某些动作在视觉上相似，但声音差异很大。
\end{itemize}

\subsection{视觉引导的音频分离}

一个很有代表性的任务是 visually guided audio source separation。直觉很简单：如果视频里能看到谁在动、谁在发声，那就可以借助视觉把混合音频拆开。

\begin{knowledgebox}{音视频分离为什么值得做}
\begin{itemize}
    \item 它把“看见谁在动”转化成“知道声音属于谁”。
    \item 它能改善会议、访谈、音乐和机器人交互中的感知质量。
    \item 它说明视频理解不仅是识别，还可以是信号分解和重建。
\end{itemize}
\end{knowledgebox}

\subsection{Audio-visual transformers 和多模态理解}

进一步地，研究者开始把图像 patch、视频 clip 和音频谱图一起送进 Transformer 结构，做统一的多模态建模。课程里列出的方向包括 attention bottleneck、audio-adaptive recognition、audio-visual masked autoencoding 等。这里的重点已经不只是“加一条音频分支”，而是把多模态表示统一进一个可学习框架。

\begin{importantbox}{多模态的两个方向}
\begin{enumerate}
    \item \textbf{判别式}：做分类、检索、定位。
    \item \textbf{生成式}：做重建、补全、mask prediction。
\end{enumerate}
\end{importantbox}

